% Draft revision of the HRES2--H2 results and statistical discussion.
% This is a separate review file; it does not modify results.tex,
% statistical_analysis_EN.tex, or the HRES2--H2 source code.
% The numerical values are restricted to the archived HRES2--H2 outputs.

\subsection{HRES2--H2 Engineering Case Study}
\label{sec:hres2_case_study_revision}

The proposed framework was evaluated on the sizing and rule-based dispatch of a
grid-connected wind--photovoltaic--electrolyzer--battery (WPEB) system located
in Baotou, Inner Mongolia, China ($41.70^\circ\mathrm{N}$,
$110.43^\circ\mathrm{E}$). The implementation optimizes the four-dimensional
mixed decision vector
\begin{equation}
    \mathbf{x} = [P_{WT}, N_{el}, P_{bat}, \tau_{bat}]^{T},
\end{equation}
where the total wind--PV capacity is fixed at $200$ MW, the electrolyzer is
represented by integer modules of $5$ MW, the battery power is discretized in
steps of $5$ MW, and the candidate battery durations are $\{1,2,4\}$ hours.
The comparison includes ten standalone metaheuristics: PSO, GWO, WOA, EHO,
ACO, ABC, ILS, SA, TS, and VNS.

For consistency with the implementation, the photovoltaic model uses a cell
temperature coefficient of $\gamma_T=-0.50\%/{}^\circ\mathrm{C}$, a nominal
operating cell temperature of $47^\circ\mathrm{C}$, and a derating factor of
$0.90$. These values replace the inconsistent $-0.45\%/{}^\circ\mathrm{C}$
value currently appearing in the companion methods text.

In the current implementation, hydrogen production is calculated from the
annual electrolyzer load. A minimum hydrogen-demand target, a hydrogen storage
tank, and a fuel-cell conversion subsystem are not modelled. Accordingly, the
case study should be interpreted as WPEB capacity sizing with annual green
hydrogen production, rather than as a complete hydrogen-storage microgrid.

\subsubsection{Experimental protocol}

Each algorithm was executed in 31 stochastic optimization runs with a maximum
budget of 1,000 iterations. Every objective evaluation simulates 8,760 hourly
dispatch steps. The archived benchmark uses one fixed synthetic 8,760-hour
weather profile for all runs; consequently, the 31 repetitions quantify
optimization stochasticity on a common scenario and do not constitute a
multi-year meteorological Monte Carlo study.

The statistical protocol comprises Shapiro--Wilk normality testing, the Friedman omnibus test, and pairwise Wilcoxon signed-rank tests. The values reported below are the unadjusted Wilcoxon $p$-values generated by the current implementation; no Holm correction is applied. The corresponding executions were treated as paired runs for the inferential comparison. The exported Mann--Whitney results provide an independent-sample sensitivity check.

\begin{table}[htbp]
\centering
\scriptsize
\caption{Archived HRES2--H2 performance summary for the proposed variants over
31 runs. AGSR values are reported as percentages.}
\label{tab:hres2_revision_summary}
\begin{tabular}{lcccccc}
\toprule
Variant & Mean LCOE & Best LCOE & Std. Dev. & Mean LCOH & AGSR range & Mean switches \\
 & (CNY/kWh) & (CNY/kWh) & (CNY/kWh) & (CNY/kg) & (\%) & \\
\midrule
DTW  & 0.267160 & 0.267159 & 0.000001 & 17.6314 & 19.9971--20.0000 & 21.00 \\
DDTW & 0.267160 & 0.267159 & 0.000005 & 17.6313 & 19.9811--20.0000 & 21.68 \\
\bottomrule
\end{tabular}
\end{table}

The archived runs support a stable low-LCOE solution for both variants. The
best recorded configuration in both cases is $P_{WT}=174.473154$ MW,
$P_{PV}=25.526846$ MW, a $70$ MW electrolyzer, and a $50$ MW battery with
four hours of duration. This configuration produces approximately
$8.862879\times10^{6}$ kg of hydrogen per year, with LCOE
$0.267159$ CNY/kWh and LCOH $17.631389$ CNY/kg. AGSR is active at, or very
close to, its $20\%$ upper bound; therefore, the result should not be
described as having a large reliability margin.

\subsubsection{Inferential results}

The Shapiro--Wilk test rejected normality for both proposed LCOE distributions
($p<10^{-11}$). Friedman tests detected significant differences among the
eleven methods in both comparisons: $\chi_F^2=129.1662$ for DTW
($p=6.910283\times10^{-23}$) and $\chi_F^2=123.4738$ for DDTW
($p=9.966826\times10^{-22}$). The proposed variants obtained mean ranks of
1.71 and 1.95, respectively.

Using the unadjusted Wilcoxon tests, all ten standalone algorithms had
$p<0.001$ against the corresponding proposed control and a higher mean rank.
Thus, the archived values support a $10/0/0$ win/tie/loss record at the
unadjusted $\alpha=0.05$ level. The independent Mann--Whitney sensitivity
test also gives $p<0.05$ for all ten baseline comparisons in both variants.
These results should be reported as evidence for the observed LCOE advantage
on this fixed engineering scenario, without claiming multiplicity-adjusted
inference.

The direct row-wise DTW--DDTW comparison yielded no significant LCOE
difference (Wilcoxon $p=0.593$). The difference in archived means was only
$7.42\times10^{-7}$ CNY/kWh, which is below the six-decimal reporting
precision. The available evidence therefore supports comparable economic
performance of the two alignment modalities in this case; it does not
establish superiority of either DTW or DDTW.

\begin{table}[p]
\centering
\scriptsize
\caption{Archived LCOE inferential comparison for HRES2--H2. The reported
pairwise values are unadjusted Wilcoxon signed-rank $p$-values.}
\label{tab:hres2_revision_inference}
\resizebox{\textwidth}{!}{\begin{tabular}{lrrrrrr}
\toprule
& \multicolumn{3}{c}{DTW} & \multicolumn{3}{c}{DDTW} \\
Algorithm & Mean LCOE & Mean rank & Raw $p$ & Mean LCOE & Mean rank & Raw $p$ \\
\midrule
Hybrid control & 0.267160 & 1.71 & -- & 0.267160 & 1.95 & -- \\
GWO & 0.268066 & 4.00 & $1.3812\times10^{-6}$ & 0.268395 & 4.06 & $1.6196\times10^{-6}$ \\
VNS & 0.267806 & 5.19 & $9.3132\times10^{-10}$ & 0.267806 & 5.00 & $1.3812\times10^{-6}$ \\
TS  & 0.267188 & 5.35 & $4.6566\times10^{-9}$ & 0.267184 & 5.52 & $9.3132\times10^{-10}$ \\
ACO & 0.273288 & 5.58 & $1.7539\times10^{-5}$ & 0.273617 & 5.68 & $4.1782\times10^{-4}$ \\
ILS & 0.267184 & 5.58 & $9.3132\times10^{-10}$ & 0.267183 & 5.58 & $3.4552\times10^{-7}$ \\
PSO & 0.277289 & 6.00 & $3.6279\times10^{-4}$ & 0.280068 & 7.00 & $8.0432\times10^{-5}$ \\
WOA & 0.277227 & 7.26 & $1.6455\times10^{-5}$ & 0.276477 & 5.97 & $2.9334\times10^{-4}$ \\
EHO & 0.277728 & 7.84 & $6.1660\times10^{-6}$ & 0.278524 & 8.06 & $2.2878\times10^{-6}$ \\
ABC & 0.272971 & 7.84 & $1.1714\times10^{-6}$ & 0.270993 & 7.34 & $1.1714\times10^{-6}$ \\
SA  & 0.286431 & 9.65 & $9.3132\times10^{-10}$ & 0.290062 & 9.84 & $9.3132\times10^{-10}$ \\
\bottomrule
\end{tabular}}
\end{table}

The inferential table intentionally reports only LCOE for the standalone
methods. The archived standalone output does not contain per-run LCOH, AGSR,
feasibility, hydrogen production, or component-sizing records. Those metrics
should not be presented as independently auditable baseline comparisons until
the corresponding per-run files are archived.

\subsubsection{Meteorological robustness extension}

A stronger engineering validation would repeat the complete experiment over a
set of hourly weather-year profiles rather than one fixed synthetic profile.
For each selected year, the weather data should be loaded into a dataframe and
passed to \texttt{HRES2Function(df\_year=...)}, while all algorithms are run
under the same scenario and documented randomization protocol. The analysis
should then report per-year and pooled distributions of LCOE, LCOH, AGSR,
hydrogen production, feasibility, and component sizing. Until this extension
is performed, the present results should be framed as a controlled synthetic
scenario study.

% Consistency corrections applied in this draft:
% (i) ten, not eight or nine, standalone algorithms are counted;
% (ii) the implementation uses gamma_T=-0.50 percent per degree Celsius;
% (iii) no hydrogen-demand, hydrogen-tank, or fuel-cell claim is made;
% (iv) standalone claims are restricted to archived LCOE statistics.
